\chapter{Analytic and Computational Tools}\label{ch:analytic-and-computational-tools}

\section{MacLane Valuations}

In his seminal work \cite{MacLane}, S.~MacLane considered and essentially solved the following problem.
Let $v_K$ be a discrete valuation on a field $K$.
What are the extensions of $v_K$ to an (arbitrary) pseudovaluation $v$ on the polynomial ring $K[x]$?
How can we describe such pseudovaluation explicitly?
These results have been applied to the problem of constructing models of curves in J.~Rüth's thesis, \cite{Rueth}. 
In this section, we give a brief overview.
We assume $K$ to be algebraically closed.


\subsection{Basic Definitions}

\begin{definition} \label{def:pseudo-valuation}
    We call $v: K[x]: \to \hat{\RR}$ a pseudovaluation if $v(1) = 0$ and
    \begin{enumerate}[(a)]
        \item $v(fg) = v(f) + v(g)$,
        \item $v(f + g) \geq \min\lbrace v(f), v(g) \rbrace$ for all $f,g\in K[x]$.
    \end{enumerate}
\end{definition}

\begin{remark}
    Let $v$ be a pseudovaluation on $K[x]$.
    \begin{enumerate} [(i)]
        \item As a consequence of the definition we get that $\p_v \coloneqq v^{-1}(\infty) < K[x]$ is a prime ideal
            and that $v$ induces a valuation 
            \begin{align}
                \bar{v}: K[x]/\p_v \to \hat{\QQ}.
            \end{align}
            If $\p_v = (0)$, then $v=\bar{v}$ is a valuation in the usual sense.
        \item Let $L|K$ be a finite field extension \textcolor{red}{[separable?]} with minimal polynomial $\varphi \in K[x]$.
            Then we have a one-to-one correspondence between the valuations $v_L: L \to \hat{\RR}$ extending $v_K$ and 
            and the pseudovaluations $v: K[x] \to \hat{\RR}$ with $v(\varphi)=\infty$.
            It is given by the mutually inverse maps
            \begin{align}
                v_L \mapsto \left( K[x] \twoheadrightarrow L \overset{v}{\to} \hat{\RR} \right)
            \end{align}
            and
            \begin{align}
                v \mapsto \left(K[x]/\p_v \overset{\bar{v}}{\to} \hat{\RR}\right).
            \end{align}
    \end{enumerate}
\end{remark}

\begin{definition}
    A pseudovaluation $v: K[x] \to \QQ$ is called a \textit{MacLane pseudovaluation} if the following conditions hold.
    \begin{enumerate}[(a)]
        \setcounter{enumi}{2}
        \item $v|_K = v_K$,
        \item $v(x) \geq 0$,
        \item $v$ is discrete, i.e. $v(K[x]) \cap \QQ = \frac{1}{e_v}\ZZ$ for some positive integer $e_v$,
        \item if $\p_v = (0)$, then the extension $k(v)|k$ has transcendence degree one.
    \end{enumerate}
    We denote the set of all MacLane pseudovaluations of $K[x]$ by $V(K[x])^*$ and the subset of MacLane valuations of $K[x]$ by $V(K[x]) \subset V(K[x])^*$.
\end{definition}

For the valuations in $V(K[x])$, there is a partial order: we say that $v \leq v'$ if and only if $v(f) \leq v'(f)$ for all $f\in K[x]$.

\begin{example}
    The Gauß valuation $v_{\scriptscriptstyle G} \in V(K[x])^*$ is minimal with respect to that partial order.
\end{example}

\begin{definition}
    Let $v\in V(K[x])$ be a MacLane valuation and let $f,g,h \in K[x]$.
    \begin{enumerate}[(i)]
        \item We say that $f$ and $g$ are \textit{$v$-equivalent} (written as $f \sim_v g$) if $v(f-g) > v(f) = v(g)$.
        \item The polynomial $f$ is said to \textit{$v$-divide} $g$ (written as $f\, |_v\, g$) if $f \sim_v fh$ for some $h \in K[x]$.
            It is called \textit{$v$-irreducible} if $f\, |_v\, gh$ implies $f\, |_v\, g$ or $f\, |_v\, h$,
            and \textit{$v$-minimal} if $f\, |_v\, g$ implies $\deg f < \deg g$.
        \item A polynomial $\phi \in K[x]$ is called a \textit{key polynomial} for $v$ is $\phi$ is monic, integral, $v$-irreducible and $v$-minimal.
    \end{enumerate}
\end{definition}

\begin{remark}
    One immediate consequence of these definitions is that any key polynomial $\phi$ for $v$ is irreducible (as an element of $K[x]$).
\end{remark}

Let $\phi$ be a key polynomial for the a valuation $v \in V(K[x])$ and let $\lambda \in \hat{\QQ}$ with $\lambda \geq v(\phi)$.
Then we can define a MacLane pseudovaluation 
\begin{align} \label{eq:def-augmentation}
    v' = [v, v'(\phi) = \lambda] \in V(K[x])^*
\end{align}
as follows.
Any polynomial $f \in K[x]$ can be written uniquely as 
\begin{align}
    f = \sum_i f_i \phi^i
\end{align}
with $\deg f_i < \deg \phi$. Then let 
\begin{align}
    v'(f) \coloneqq \min_i(v(f_i) + i\cdot \lambda).
\end{align}

It holds that $v\leq v'$ and $v < v'$ if and only if $v(\phi) < \lambda$.

\begin{definition}
    The MacLane pseudovaluation $v'$ defined in (\refeq{eq:def-augmentation}) is called the \textit{augmentation} of $v$ with respect to $(\phi,\lambda)$.
    If $v < v'$, $v'$ is called a \textit{proper augmentation}.
\end{definition}

\begin{definition}
    A MacLane pseudovaluation $v$ is called \textit{inductive} if there exists a chain of MacLane pseudovaluations
    \begin{align}
        v_{\scriptscriptstyle G} = v_0 \leq v_1 \leq \cdots \leq v_n = v,
    \end{align}
    such that $v_i$ is a proper augmentation of $v_{i-1}$ for $i=1,\dots, n$.
    We call $v_0,\dots,v_n$ the \textit{augmentation chain} representing $v$ and write \begin{align}
        v = [v_0,v_1(\phi_1) = \lambda_1,\dots, v_n(\phi_n) = \lambda_n].
    \end{align}
    Here it is understood $v_0 = v_{\scriptscriptstyle G}$, $v_n = v$, that $\phi_i$ is a key polynomial for $v_{i-1}$, $\lambda_i > v_{i-1}(\phi_i)$
    and that $v_i = [v_{i-1}, v_i(\phi_i) = \lambda_i]$. 
    Also, $\lambda_i \neq \infty$ for $i < n$.
    We say that the augmentation chain $v_0,\dots,v_n$ is \textit{minimal} if
    \begin{align}
        \deg \phi_i < \cdots < \deg \phi_n.
    \end{align}
\end{definition}

\begin{theorem}[{\cite[Theorem 8.1 and Theorem 15.3]{MacLane}}]
    Let $v \in V(K[x])^*$.
    Then there exists a unique minimal augmentation chain representing $v$.
    In particular, $v$ is inductive.
\end{theorem}


\subsection{Discs, Discoids and Inductive Valuations}

\begin{definition}
    \begin{enumerate}[(i)]
        \item Let $\alpha \in K$ and let $\lambda \in \hat{\QQ}$. 
            Then
                \begin{align}
                    D(\alpha, \lambda) \coloneqq \lbrace x \in K:\, v_K(x - \alpha) \geq \lambda\rbrace
                \end{align}
            is the \textit{(closed) disc} with centre $\alpha$ and radius $\lambda$.
        \item Let $\phi \in K[x]$ be a monic irreducible polynomial and let $\lambda \in \hat{\QQ}$.
            The set 
            \begin{align}
                D(\phi, \lambda) \coloneqq \lbrace x \in K:\, v_K(\phi(x)) \geq \lambda \rbrace
            \end{align}
            is called the \textit{discoid} with centre $\phi$ and radius $\lambda$.
    \end{enumerate}
\end{definition}

\begin{remark}
    Any disc with centre in $K$ is a discoid.
    Coversely, if $\phi = x - \alpha$ then $D(\phi, \lambda)$ coincides with the disc $D(\alpha, \lambda)$.
    For arbitrary non-constant $\phi$, we have the following generalisation:
    Let $\alpha_1,\dots,\alpha_n$ be the (pairwise distinct) roots of $\phi \in K[x]$
    and let $\mu_1,\dots,\mu_n \in \hat{\QQ}$ be minimal\footnote{
        Our definition of discs imply that discs with large radius are smaller than those with small radius.
        For instance, the disc with centre $0\in K$ and radius $0$ is rather big, as it consists of all integral elements.
    } such that $D(\alpha_i,\mu_i) \subset D(\phi,\lambda)$.
    Then
    \begin{align}
        D(\phi,\lambda) = \bigcup_{i=1}^n D(\alpha_i,\mu_i),
    \end{align}
    see \cite[Lemma 4.43]{Rueth}.
\end{remark}